\subsection{Field variation: an exact finite-character reduction}
\label{sec:kt-field-variation}

The field-selection question in this subsection is deliberately narrower
than the full automorphic exclusion question.  It asks for an imaginary
quadratic field whose class number and two specified normalized
generalized Bernoulli numbers are units.  Even a positive answer would
need a separate theorem to exclude every stable rank-three congruence.

Let $\ell>3$ be prime, let $e\geq2$ be even, and assume the actual
\texttt{GD} congruence condition
\[
                         (\ell-1)\nmid 2e.                 \tag{FV.1}
\]
Let $D<0$ be a fundamental discriminant, $f=|D|$, with $\ell\nmid f$,
and put $\chi=\chi_D$.  Fix the Teichm\"uller character
$\omega:(\mathbf Z/\ell\mathbf Z)^\times\to\mathbf Z_\ell^\times$.
Congruences involving its values use this fixed $\ell$-adic embedding.
Write $\beta_{n,\chi}=B_{n,\chi}/n$ for $n\geq1$.

\begin{proposition}[KT-FIELD-CHARACTER-REDUCTION]
For every odd $n\geq3$ with $n\not\equiv1\pmod{\ell-1}$, let $k$ be
the unique odd integer in $[3,\ell-2]$ congruent to $n$ modulo
$\ell-1$.  Then
\[
 \beta_{n,\chi}\equiv\beta_{k,\chi}
       \equiv B_{1,\chi\omega^{k-1}}\pmod\ell.              \tag{FV.2}
\]
All three quantities are $\ell$-integral.  The character in the last
term is primitive of conductor $f\ell$.  In particular,
\[
 \beta_{e+1,\chi_D}\equiv B_{1,\chi_D\omega^e},\qquad
 \beta_{2e+1,\chi_D}\equiv B_{1,\chi_D\omega^{2e}}
                                      \pmod\ell.           \tag{FV.3}
\]
There is no hypothesis comparing $\ell$ with the original weight $n$.
\end{proposition}

\begin{proof}
We use Sun's Lemma~8.1(a), with its difference order equal to one:
if $b\geq1$ and $(\ell-1)\nmid b$, then
\[
 (1-\chi(\ell)\ell^{b-1})\beta_{b,\chi}
 \equiv
 (1-\chi(\ell)\ell^{b+\ell-2})\beta_{b+\ell-1,\chi}
                                  \pmod\ell.               \tag{FV.4}
\]
The conductor hypothesis is exactly $\ell\nmid f$.  The published
statement is in Z.-H.~Sun, \emph{Congruences concerning Bernoulli
numbers and Bernoulli polynomials}, Discrete Applied Mathematics
105 (2000), 193--223, pp.~219--220.  Its full proof was already
checked for KT-TWOFIELD-FAMILY; the exact statement and proof are
reread for the present application.

First all $\beta_{m,\chi}$ for odd $m\geq1$ are $\ell$-integral.
For odd $1\leq b\leq\ell-2$, this follows from
\[
 B_{b,\chi}=f^{b-1}\sum_{a=1}^{f}\chi(a)B_b(a/f).
\]
Indeed the ordinary Bernoulli numbers $B_j$ for $0\leq j\leq\ell-2$
are $\ell$-integral: their elementary recurrence only divides by
$j+1<\ell$.  Also $b$ and $f$ are units.  Induction along each odd
residue class using (FV.4) proves the assertion for every odd $m$.
At the first step from $b=1$, the factor $1-\chi(\ell)$ can be zero;
this causes no difficulty because its product with $\beta_{1,\chi}$
is integral, while the Euler factor at $b+\ell-1$ is a unit.
At every later step both Euler factors are units congruent to one.
Since $\chi$ is odd, $B_{m,\chi}=0$ for even $m$, including $m=0$.
Consequently every $B_{m,\chi}$ is $\ell$-integral.

For $n$ and $k$ in the proposition, (FV.4) can be applied from $k$
to $n$ using only indices at least three.  Both Euler factors are
then congruent to one; hence $\beta_{n,\chi}\equiv\beta_{k,\chi}$.
This includes indices divisible by $\ell$: no division by their
residue classes is performed.

Put
\[
                        r=1+\ell(k-1),\qquad M=f\ell.
\]
Then $r\equiv k\pmod{\ell-1}$, $r\equiv1\pmod\ell$, and $r\geq3$.
The same congruence gives $\beta_{r,\chi}\equiv\beta_{k,\chi}$.
The generalized Bernoulli generating series, summed over $\ell$
complete periods of $\chi$, gives the exact identity
\[
 \frac1M\sum_{a=1}^{M}\chi(a)a^r
   =\sum_{j=0}^{r}\binom rj B_{r-j,\chi}\frac{M^j}{j+1}.
                                                               \tag{FV.5}
\]
For $j\geq1$ one has $j-v_\ell(j+1)\geq1$.  The integrality just
proved therefore implies
\[
 \frac1M\sum_{a=1}^{M}\chi(a)a^r\equiv B_{r,\chi}\pmod\ell.
                                                               \tag{FV.6}
\]
Removing multiples of $\ell$ from this sum changes its divided
value by
$\chi(\ell)\ell^{r-1}f^{-1}\sum_{b=1}^{f}\chi(b)b^r$,
which is divisible by $\ell$.

For every $\ell$-adic unit $a$,
$\omega(a)\equiv a^\ell\pmod{\ell^2}$.  Indeed write
$a=\omega(a)(1+\ell z)$ and raise to the $\ell$th power.
The primitive character $\theta=\chi\omega^{k-1}$ has conductor
$f\ell$: the two factors have coprime conductors, and $k-1$ is
nonzero modulo $\ell-1$.  Its defining formula is
\[
 B_{1,\theta}=\frac1{f\ell}
       \sum_{\substack{1\leq a\leq f\ell\\\ell\nmid a}}
               \chi(a)\omega(a)^{k-1}a.
\]
Replacing $\omega(a)$ by $a^\ell$ changes each summand by a
multiple of $\ell^2$, hence changes the divided sum by a multiple
of $\ell$.  Combining this observation with (FV.6) proves
\[
 B_{1,\theta}\equiv B_{r,\chi}
       =r\beta_{r,\chi}\equiv\beta_{k,\chi}\pmod\ell.
\]
It also proves the asserted integrality of $B_{1,\theta}$.
Finally (FV.1) excludes both $e\equiv0$ and $2e\equiv0$
modulo $\ell-1$, so the proposition applies to $e+1$ and $2e+1$.
\end{proof}

\begin{corollary}[The exact simultaneous field-selection requirement]
Let $\mathcal D$ be any prescribed family of negative fundamental
discriminants prime to $\ell$, with any additional local splitting
or parity conditions.  A member $D\in\mathcal D$ satisfies
\[
 \ell\nmid h(D),\qquad
 \beta_{e+1,\chi_D},\ \beta_{2e+1,\chi_D}\in\mathbf Z_\ell^\times
                                                               \tag{FV.7}
\]
if and only if
\[
 B_{1,\chi_D}\,B_{1,\chi_D\omega^e}\,
 B_{1,\chi_D\omega^{2e}}\not\equiv0\pmod\ell.               \tag{FV.8}
\]
The first twisted character $\chi_D\omega^e$ is never quadratic
under (FV.1).
\end{corollary}
\begin{proof}
The imaginary quadratic class-number formula is
$B_{1,\chi_D}=-2h(D)/w(D)$, with $w(D)\in\{2,4,6\}$.
Since $\ell>3$, its unit condition is exactly $\ell\nmid h(D)$.
Now apply (FV.3).  On the $\ell$-part of the conductor the order
of $\chi_D\omega^e$ is
$(\ell-1)/\gcd(e,\ell-1)>2$, by (FV.1), proving the last assertion.
\end{proof}

\paragraph{What the checked field-variation theorems supply.}
Wiles, \emph{On class groups of imaginary quadratic fields}, JLMS
92 (2015), 411--426, Theorem~A (accepted manuscript
Theorem~0.0.1, pp.~1--2), supplies the first factor in (FV.8),
with specified finite sets of odd split, inert, and ramified primes
subject to its three explicit residue restrictions.  Its original
statement allows $\ell$ in a local set; in particular one can ask
that $\ell$ split.  It is a quadratic-character theorem and does
not provide either nonquadratic twist in (FV.8).  Beckwith's
quantitative refinement, Theorem~1.1 and Corollary~1.2, concerns
the same class-number factor; its restatement excludes $\ell$
from the specified local sets.

Bruinier, Duke Math.~J.~98 (1999), 595--611, Theorem~3
(author PDF p.~11), treats one integral primitive half-integral
weight Hecke eigenform at a time.  For a Cohen series of weight
$k+\tfrac12$ and an odd prescribed local prime $p$, its exceptional
set condition is
\[
 \ell\nmid p(p-1)(p^{2k}-1)(p^{2k-2}-1).                    \tag{FV.9}
\]
This is obtained by substituting $1+p^{2k-1}$ for the Hecke
eigenvalue and factoring the two possible eigenvalue differences.
The theorem also requires the given series to be integral and
primitive at $\ell$; rescaling by a nonunit cannot establish
unit values in the original normalization.  Its Dirichlet
specialization Theorem~6 is stated for even $k$ and real fields,
with an explicit unit seed.  The odd-weight analogue is mentioned
afterwards, rather than supplied as an unconditional imaginary
field theorem.  These precise conditions are recorded in the
separate source audit.

The limitation (FV.9) is genuine for this application of the
theorem: $\ell=17$, $e=4$ satisfies (FV.1), but for the second
index $k=2e+1=9$, every $p\ne17$ has $p^{2k-2}=p^{16}\equiv1$.
Thus this version of the prescribed-local-data theorem cannot
be applied, even though both weights are smaller than $\ell$.
This does not assert that the desired fields fail to exist.

Burungale--Sun, Math.~Ann.~378 (2020), 317--358,
Theorem~5.2 and Corollary~5.4, give quantitative nonvanishing as
\emph{all} characters of a growing conductor vary.  Their family
does not restrict to quadratic characters; for prime conductor
there is only one nontrivial quadratic character, so their lower
bound cannot force that character to contribute.  Their
fixed-base-character prime-power tower results likewise vary
characters of growing nonquadratic order.  Neither statement is
a simultaneous assertion for the three factors in (FV.8).
There is also a direct conductor obstruction to applying
Theorem~5.2 to the fixed character $\omega^e$: that theorem
requires $\ell\nmid N$ for the fixed period $N$, whereas
$\omega^e$ has conductor $\ell$ under (FV.1).

\paragraph{The remaining precise implication.}
The reviewed sources do not establish existence in (FV.8) for
each admissible residue class $e\bmod(\ell-1)$ with the required
local conditions.  Separate infinite sets of discriminants with
individual unit values do not imply that their intersection is
nonempty.  Multiplying two Cohen $q$-series forms a convolution
of coefficients, not the same-discriminant product required in
(FV.8).  This is the exact additional field-selection theorem
needed by this Bernoulli strategy.  It is a scope statement about
the checked results, not a claim that no stronger theorem exists.

Fretwell--Roberts Theorem~3.6 supplies the direction from
Bernoulli divisibility to an Eisenstein congruence. Their
necessary-condition Theorem~3.5 has a large-prime hypothesis
relative to the actual weight; Kummer reduction alone does not
transport a high-weight system into that range. The missing
ordinary implication has now been proved by a separate argument:
KT-ORDINARY-RANK2-UNIT-EXCLUSION uses the exact auxiliary strict
group and an opposite stable lattice, with no weight bound.
KT-AUXILIARY-UNIT-NO-TYPEII matches it to actual Type-II source
points, under GD, $N=2$ and $\psi=1$, using the $a=2e$ unit.
The $a=e$ version also excludes the specified ordinary rank-two
seeds of the symmetric-square construction. Uniform field
selection with the actual GD restrictions, and exclusion of
the remaining stable rank-three sources, are still required.

\begin{lemma}[KT-FIELD-SPECIALIZATION-LIMIT]
Avoidance of common height-one primes by two power series does
not imply that one of their values is a unit at every integral
specialization, even if the two series use independent variables.
\end{lemma}
\begin{proof}
In $A=\mathbf Z_\ell[[X,Y]]$ take $f=\ell+X$ and $g=\ell+Y$.
The quotient $A/(f)$ is $\mathbf Z_\ell[[Y]]$ by the continuous
substitution $X=-\ell$; hence $(f)$ is a nonzero height-one
prime. Any height-one prime containing $f$ must equal $(f)$,
whereas the image $\ell+Y$ of $g$ in that quotient is nonzero.
Thus no height-one prime contains both. At the continuous
integral specialization $X=Y=0$ both values equal $\ell$,
so both are nonunits.
\end{proof}
